%% kapa: exam=1 round=29 year=2018 period=1 form=A subject=economics subject_ko=경제학원론 seq=22 number=62
%% source: Q-Net 감정평가사 기출문제 (cst00309) artlSeq=5202808 file=extracted/1008/5202808/2018년 제29회 감정평가사 1차시험 기출문제/2018년 제29회 감정평가사 1차시험 기출문제/1교시 A형.pdf
%% conversion: kordoc 4.24.51 -> tools/gen1.py (auto); review=proofread-sonnet; flags=none
\begin{question}{62}
물가지수에 관한 설명으로 옳은 것은?
\begin{choices}
  \item GDP 디플레이터에는 국내산 최종 소비재만이 포함된다.
  \item GDP 디플레이터 작성 시 재화와 서비스의 가격에 적용되는 가중치가 매년 달라진다.
  \item 소비자물가지수 산정에는 국내에서 생산되는 재화만 포함된다.
  \item 소비자물가지수에는 국민이 구매한 모든 재화와 서비스가 포함된다.
  \item 생산자물가지수에는 기업이 구매하는 품목 중 원자재를 제외한 품목이 포함된다.
\end{choices}
\end{question}
\answer{2}
